\begin{textsample}{NEG Dim 6 – ro – Score -2.00 – ro/ro\_ef\_78.txt}  \label{ex:f6_neg_151}
9.5.9 Conteúdo básico comum da Educação Física A Educação Física será composta de seis unidades temáticas de acordo com o que é proposto na BNCC.
Tais unidades temáticas e seus objetos de conhecimentos estão apresentados conforme quadro abaixo, que correspondem aos temas que devem ser abordados nas aulas de Educação Física no Estado de Rondônia no ensino do 1º ao 9º anos do Ensino Fundamental.
Brincadeiras e jogos : Essa unidade explora as atividades voluntárias exercidas dentro dos limites de tempo e espaço, caracterizadas pela criação e alteração das regras.
Trata -se de práticas corporais populares, motivadoras e desafiadoras, proporcionando ao aluno processos criativos significativos e modificadores no imaginário, na realidade e no presente.
Deve -se observar a diferença entre o jogo como conteúdo programático e a ferramenta auxiliar de ensino.
A BNCC contempla essa temática como conteúdo de extremo valor e relevância no aprendizado da cultura corporal.
Esportes : Trata -se da manifestação corporal mais difundida em todo o mundo, porém ela institucionaliza temas lúdicos da cultura corporal que envolvem códigos, significados e regras pré-estabelecidas.
O esporte é um fenômeno plural que pode ocorrer em diversos contextos de prática, em níveis diferenciados de exigência.
Portanto, a grande variabilidade de sentidos correlacionados com o esporte indicou a necessidade de classificá -lo de modo mais detalhado em : Esporte de Rendimento, Esporte educacional e esporte de participação.
Para a BNCC, o esporte segue o modelo que foi baseado nos critérios de cooperação, interação com o adversário, desempenho motor e objetivos táticos da ação, privilegiando -se as ações motora intrínsecas do aluno aliadas às exigências motrizes semelhantes na prática dos esportes, em detrimento ao esporte que visa performance e resultados otimizados.
Para facilitar a compreensão do que caracteriza cada uma das modalidades, estas se apresentam divididas em 07 categorias : Marca : conjunto de modalidades que se caracterizam por comparar os resultados registrados em segundos, metros ou quilos ( patinação de velocidade, todas as provas do atletismo, remo, ciclismo, levantamento de peso etc. ).
Precisão : conjunto de modalidades que se caracterizam por arremessar/lançar um objeto, procurando acertar um alvo específico, estático ou em movimento, comparando -se o número de tentativas empreendidas, a pontuação estabelecida em cada tentativa ( maior ou menor do que a do adversário ) ou a proximidade do objeto arremessado ao alvo ( mais perto ou mais longe do que o adversário conseguiu deixar ), como nos seguintes casos : bocha, curling, golfe, tiro com arco, tiro esportivo etc.
Técnico-combinatório : reúne modalidades nas quais o resultado da ação motora comparado é a qualidade do movimento segundo padrões técnico-combinatórios ( ginástica artística, ginástica rítmica, nado sincronizado, patinação artística, saltos ornamentais etc. ).
Rede/quadra dividida ou parede de rebote : reúne modalidades que se caracterizam por arremessar, lançar ou rebater a bola em direção a setores da quadra adversária nos quais o rival seja incapaz de devolvê -la da mesma forma ou que leve o adversário a cometer um erro dentro do período de tempo em que o objeto do jogo está em movimento.
Alguns exemplos de esportes de rede são voleibol, vôlei de praia, tênis de campo, tênis de mesa, badminton e peteca.
Já os esportes de parede incluem pelota basca, raquetebol, squash etc.
Campo e taco : categoria que reúne as modalidades que se caracterizam por rebater a bola lançada pelo adversário o mais longe possível, para tentar percorrer o maior número de vezes as bases ou a maior distância possível entre as bases, enquanto os defensores não \textbf{recuperam} o controle da bola, e, assim, somar pontos ( beisebol, críquete, softbol etc. ).
Invasão ou territorial : conjunto de modalidades que se caracterizam por comparar a capacidade de uma equipe introduzir ou levar uma bola ( ou outro objeto ) a uma meta ou setor da quadra/campo defendida pelos adversários ( gol, cesta, touchdown etc. ), protegendo, simultaneamente, o próprio alvo, meta ou setor do campo ( basquetebol, frisbee, futebol, futsal, futebol americano, handebol, hóquei sobre grama, polo aquático, rúgbi etc. ).
Combate : reúne modalidades caracterizadas como disputas nas quais o oponente deve ser subjugado, com técnicas, táticas e estratégias de desequilíbrio, contusão, imobilização ou exclusão de um determinado espaço, por meio de combinações de ações de ataque e defesa ( judô, boxe, esgrima, taekwondo etc. ).
Ginásticas : está constituída na história humana desde a pré-história, mantendo -se na Idade Média, fundamentando -se na Idade Moderna e perpassando por sistematizações na Idade Contemporânea.
Trata -se de uma prática corporal que exercita o corpo e visa à exploração das possibilidades acrobáticas e expressivas do corpo.
A BNCC denota que as ginásticas de competição ( artística, rítmica, acrobática e aeróbica de trampolim ) são classificadas como esportes técnicos-combinatórios.
Essa temática subdivide -se como conteúdo pedagógico em : Ginástica geral : também conhecida como ginástica para todos, reúne as práticas corporais que têm como elemento organizador a exploração das possibilidades acrobáticas e expressivas do corpo, a interação social, o compartilhamento do aprendizado e a não competitividade.
Podem ser constituídas de exercícios no solo, no ar ( saltos ), em aparelhos ( trapézio, corda, fita elástica ), de maneira individual ou coletiva, e combinam um conjunto bem variado de piruetas, rolamentos, paradas de mão, pontes, pirâmides humanas etc.
Integram também essa prática os denominados jogos de malabar ou malabarismo.
Ginástica de condicionamento físico : se caracteriza pela exercitação corporal orientada à melhoria do rendimento, à aquisição e à manutenção da condição física individual ou à modificação da composição corporal.
Geralmente, é organizada em sessões planejadas de movimentos repetidos, com frequência e intensidade definidas.
Pode ser orientada de acordo com uma população específica, como a ginástica para gestantes, ou atrelada a situações ambientais determinadas, como a ginástica laboral.
Ginásticas de conscientização corporal : reúnem práticas que empregam movimentos suaves e lentos, tal como a recorrência a posturas ou à conscientização de exercícios respiratórios, voltados para a obtenção de uma melhor percepção sobre o próprio corpo.
Algumas dessas práticas que constituem esse grupo têm origem em práticas corporais milenares da cultura oriental.
Danças : manifestação cultural que reúne movimento corporal, música, ritmo, expressão corporal e uma variedade de sentimentos para quem exerce sua prática, é um patrimônio não material que irradia arte e beleza.
Também pode ser compreendida como uma expressão de movimentos corporais organizados, partindo de experiências significativas do movimento humano.
As danças podem ser realizadas de forma individual, em duplas ou em grupos, sendo essas duas últimas as formas mais comuns.
Diferentes de outras práticas corporais rítmico-expressivas, elas se desenvolvem em \textbf{codificações} particulares, historicamente \textbf{constituídas}, que permitem identificar movimentos e ritmos musicais peculiares associados a cada uma delas. ( BNCC, 2017 ).
Lutas : Trata -se de disputas corporais, nas quais os participantes empregam técnicas, táticas e estratégias específicas para imobilizar, desequilibrar, atingir ou excluir o oponente de um determinado espaço, combinando ações de ataque e defesa dirigidas ao corpo do adversário.
Dessa forma, além das lutas presentes no contexto comunitário e regional, podem ser tratadas lutas brasileiras ( capoeira, huka-huka, luta marajoara etc. ), bem como lutas de diversos países do mundo ( judô, aikido, jiu-jítsu, muay thai, boxe, Chinese boxing, esgrima, kendo etc. ).
Práticas Corporais de Aventura : São expressões e formas de experimentação corporal centradas nas perícias e proezas provocadas pelas situações de imprevisibilidade que se apresentam quando o praticante interage com um ambiente desafiador.
Para Cantorani e Pilatti ( 2005 ), ao longo da evolução humana, sempre existiram práticas que envolviam desafio e aventuras em que fortes emoções se faziam presentes.
Na BNCC, optou -se por diferenciá -las com base no ambiente de que necessitam para ser realizadas : na natureza e urbanas.
As práticas de aventura na natureza se caracterizam por explorar as incertezas que o ambiente físico cria para o praticante na geração da vertigem e do risco controlado, como em corrida orientada, corrida de aventura, corridas de mountain bike, rapel, tirolesa, arborismo etc.
Já as práticas de aventura urbanas exploram a “ paisagem de cimento ” para produzir essas condições ( vertigem e risco controlado ) durante a prática de parkour, skate, patins, bike etc.

% matched lemmas: codificação, constituído, recuperar
\end{textsample}
